\clearpage
\chapter*{Erratum to Exposé no. 212}
\addcontentsline{toc}{chapter}{Erratum to Exposé no. 212}
\setcounter{footnote}{2}
\src{301}

\noindent\textbf{5. TDTE III: Quotient preschemes.}

\begin{center}
[Bourbaki Seminar, vol. 13, 1960/61, no. 212, 20 pp.]
\end{center}

\noindent\textbf{Page 15, assertion (i).} --- The fact that $Y=X/R$ is
quasi-projective over $S$ appears, for the present, to have been proved only
when $R$ arises from an equivalence relation.

\noindent\textbf{Pages 18 and 19.} --- As we point out at the end of the
following exposé, the conjecture under consideration is definitely false. The
``positive fact'' noted in Remark 8.1 appears to have been proved
simultaneously by various authors (Nagata, Rosenlicht, Grothendieck, \ldots).

Let us point out that the application in TDTE V of the theory of passage to
the quotient developed here to the construction of Picard schemes may also be
replaced by an appropriate use of Hilbert schemes (cf. SHMT). As was stated in
§ 8, the most important gap in the theory presented
\src{302}
here is the lack of a criterion for the existence of quotients by a nonproper
equivalence relation, such as equivalence relations arising from certain
actions of the projective group. An important theorem in this direction has
been obtained by Mumford\footnote{Mumford (David). --- An elementary theorem
in geometric invariant theory, \emph{Bull. Amer. Math. Soc.}, vol. 67, 1961,
pp. 483--486.}. For a refinement of his result and various applications to the
theory, see SHMT.
